\documentclass[12pt,a4paper,fontset=fandol]{ctexart}
\usepackage{geometry}
\geometry{a4paper,left=2.5cm,right=2.5cm,top=2.5cm,bottom=2.5cm}
\usepackage{hyperref}
\usepackage{enumitem}
\hypersetup{pdfauthor={},pdftitle={AI工具使用详情}}
\pagestyle{plain}
\setlength{\parindent}{2em}
\linespread{1.18}
\begin{document}
\begin{center}
{\LARGE\bfseries AI工具使用详情\par}
\end{center}

\section*{1. 所用 AI 工具名称、版本或型号}

\begin{itemize}[leftmargin=2em]
\item Codex（OpenAI），模型版本 5.6 sol，使用日期 2026 年 9 月 10 日至 9 月 13 日。
\item DeepSeek V4.1 Flash（深度求索），使用日期 2026 年 9 月 13 日。
\end{itemize}

\section*{2. 具体使用目的和环节}

本队由参赛队员主导全部建模工作：题目分析、模型选择、假设与原理推导、算法设计、参数设定、结算与收费口径、结果解释及结论均由队员完成。AI 工具仅作为辅助手段，承担以下辅助性、机械性工作：

\begin{enumerate}[leftmargin=2em]
\item 程序实现：将队员已经确定的算法思路和输入输出要求转化为可运行的程序代码，包括数据读取与时间标签解析、预测、滚动优化、实际反馈与结算、结果导出等程序。
\item 工程优化：按队员逐项指定的工作包，完成计算加速、结构缓存、求解器接口等机械性改动，并核对改动前后结果是否一致。
\item 文字辅助：按队员指定的章节范围和写作要求，协助整理文字初稿，供队员修改与定稿。
\item 图表与排版：依据队员已有的计算结果绘制图件，并完成论文排版。
\end{enumerate}

AI 不参与模型选择与算法设计，不作出任何决策，也不更改任何已确定的方案。

\section*{3. 主要提示方式与使用过程说明}

队员以自然语言给出任务，通常包含三部分：任务范围、必须遵守的边界、以及可读取的本地材料位置。AI 读取现有稿件、已确定的模型口径与已有的计算结果后产出可编辑的程序或文字稿，再由队员检查。边界条件由队员显式声明，例如：不实现未经确定的方案、原始附件只读、不得改动已确定的结论。论文风格方面，队员提供往年优秀论文作为体裁与图表风格参考，该参考仅用于篇章组织、行文层次与图表样式，不移植其中的模型、结论或原文。以下为典型交互示例（节选、非逐字记录）。

\textbf{示例一（程序实现）。} 队员给出已确定的算法口径、输入输出要求与边界条件，要求据此编写程序。AI 返回实现与测试，队员在本地运行核对，对不符合预期的部分给出具体反例，AI 据此修正。

\textbf{示例二（文字辅助）。} 队员指定需要整理的章节范围与论文体裁，并提供既有的计算结果。AI 起草该章节的文字初稿，队员审阅后决定采纳、修改或退回。

\textbf{示例三（边界约束）。} 当任务涉及尚未确定的方案时，队员在提示中声明不得代替其作出决定。AI 仅整理备选方案与计算结果，不自行选定方案。

\section*{4. 对 AI 输出的采纳、人工修改和核验的主要情况}

本队对 AI 输出实行先核对后使用的流程，所有 AI 生成内容均不作为最终结果直接采用：

\begin{enumerate}[leftmargin=2em]
\item 程序代码由队员逐段阅读、运行并补充测试后使用；关键结论均由队员依据原题与计算结果独立复算核对。
\item 文字初稿由队员逐句审阅，模型表述、公式、数值与结论均以队员确定的口径为准，由队员修改定稿。按官方要求，语言润色部分不在此列。
\item 图表由队员核对所依据的数据来源与数值后采用。
\end{enumerate}

队员重点核对的内容包括：时间标签与区间映射、效率与单位口径、预测的因果可用时点、储能与末态约束、各项费用分项，以及摘要与结论中的数值陈述。

\end{document}
